\section{Síntese crítica da literatura e posicionamento da dissertação}
\label{sec_sintese_critica_posicionamento_dissertacao}

A revisão da literatura apresentada ao longo deste capítulo evidencia que os estudos sobre operações de \textit{rendezvous} e \textit{berthing} evoluíram de formulações mais clássicas de dinâmica relativa para abordagens mais amplas, envolvendo percepção relativa, controle de atitude, planejamento de trajetória, autonomia e interação robótica. As referências revisadas indicam que a literatura avançou tanto na formulação matemática do movimento relativo entre espaçonaves quanto na modelagem e no controle de sistemas robóticos espaciais acoplados.

No campo da dinâmica relativa, os trabalhos clássicos estabeleceram as bases matemáticas para a descrição do movimento entre perseguidor e alvo. Em estudos posteriores, o problema passou a incorporar restrições operacionais, requisitos de segurança e condições associadas à aproximação de alvos cooperativos e não cooperativos. Na área de percepção relativa, a literatura passou a apresentar métodos voltados à estimação de posição, atitude, velocidade relativa e, em alguns casos, da própria geometria do alvo, utilizando sensores passivos e ativos, além de técnicas mais recentes baseadas em aprendizado de máquina.

Além disso, a literatura revisada mostra que as operações de \textit{rendezvous} e \textit{berthing} vêm sendo tratadas em contextos operacionais mais amplos, envolvendo diferentes níveis de cooperação do alvo, troca de informações entre sistemas embarcados, requisitos de segurança e maior autonomia decisória ao longo da missão.
Embora nem todas essas abordagens constituam o núcleo da metodologia adotada nesta dissertação, sua inclusão na revisão permite situar o estado da arte da área e delimitar com maior clareza o recorte adotado no presente trabalho.

No que se refere à robótica espacial, a literatura revisada evidencia que os sistemas espaçonave--manipulador devem ser tratados como sistemas dinamicamente acoplados, nos quais o movimento das juntas produz reações sobre a base e altera a dinâmica global do sistema. As referências analisadas indicam ainda que estratégias como \textit{reaction null space}, controle \textit{reactionless} e abordagens adaptativas ampliaram a capacidade de tratar problemas de captura, pós-captura e interação com alvos desconhecidos, preservando, tanto quanto possível, o comportamento dinâmico da base.

Entretanto, a revisão também mostra que parte significativa da literatura trata esses problemas de forma parcial. Em muitos trabalhos, a dinâmica relativa orbital é formulada separadamente da dinâmica de atitude da espaçonave perseguidora. Em outros, a modelagem do manipulador robótico é desenvolvida sem integração explícita com o contexto de aproximação orbital de curta distância. Há ainda referências voltadas à percepção relativa, à estimação de pose ou ao planejamento autônomo que, embora importantes, não tratam de forma integrada os efeitos dinâmicos do sistema completo quando a espaçonave perseguidora transporta ou opera um manipulador robótico em ambiente orbital.

Dessa forma, a principal lacuna observada na literatura revisada não está na ausência de estudos sobre os diferentes elementos do problema, mas na limitada integração entre eles dentro de uma mesma formulação de modelagem e análise. Em particular, permanece relevante o estudo de formulações que articulem a dinâmica relativa entre perseguidor e alvo, a dinâmica de atitude da espaçonave perseguidora no referencial orbital e a dinâmica acoplada de um manipulador robótico embarcado, com vistas à compreensão do comportamento global do sistema durante operações de aproximação e interação com o alvo.

É nesse contexto que se insere a presente dissertação. O trabalho concentra-se na modelagem integrada do sistema formado por uma espaçonave perseguidora equipada com manipulador robótico, considerando o contexto de operações de \textit{rendezvous} e \textit{berthing} em ambiente orbital. O foco recai sobre a articulação entre dinâmica relativa, dinâmica de atitude e modelagem do manipulador, de modo a representar de forma consistente os acoplamentos envolvidos nas fases de aproximação e nas condições que antecedem a interação com o alvo.

Assim, a contribuição principal desta dissertação não está na proposição de uma arquitetura completa de autonomia, nem no desenvolvimento de um sistema embarcado de percepção relativa ou de tomada de decisão em tempo real. O objetivo principal está na construção e na análise de um modelo dinâmico coerente com o problema físico estudado, capaz de servir de base para investigações posteriores em controle, planejamento e operação robótica espacial. Dessa forma, o presente trabalho adota um escopo mais delimitado, concentrando-se na formulação matemática e na análise do sistema acoplado em ambiente orbital.

De modo geral, a revisão da literatura evidencia que, embora a área tenha avançado de forma significativa em dinâmica relativa, percepção, controle e robótica espacial, ainda permanece relevante o desenvolvimento de abordagens que privilegiem a integração entre esses domínios. Nesse sentido, a presente dissertação busca contribuir ao tratar de forma articulada a dinâmica orbital relativa, a atitude da espaçonave perseguidora e a dinâmica do manipulador robótico, definindo, assim, seu posicionamento no conjunto das pesquisas dedicadas às operações de \textit{rendezvous} e \textit{berthing} com apoio de sistemas robóticos espaciais.